\documentclass[12pt,a4paper]{article}
\usepackage[utf8]{inputenc}
\usepackage[T5]{fontenc}
\usepackage[vietnamese]{babel}
\usepackage{amsmath, amssymb, amsfonts}
\usepackage{enumitem}

\newcommand{\R}{\mathbb{R}}
\newcommand{\N}{\mathbb{N}}
\newcommand{\ip}[2]{\left\langle #1,#2\right\rangle}
\newcommand{\norm}[1]{\left\|#1\right\|}

\begin{document}

\begin{center}
\Large\textbf{BÀI GIẢI MẪU ÔN TẬP CUỐI KỲ}
\end{center}

\underline{Câu 1.} Trên $V=\R^5$, cho tập hợp
\[
W = \left\{ X = (x_1; x_2; x_3; x_4; x_5) \in \R^5 \left| \begin{matrix} x_1 - x_2 + 2x_3 - x_4 + x_5 = 0 \\ 2x_1 - x_2 + 3x_3 - 2x_4 + x_5 = 0 \end{matrix} \right. \right\}
\]
Chứng minh rằng $W$ là không gian véc tơ con của $V$. Tìm hệ sinh, cơ sở và xác định số chiều cho $W$.

\underline{Giải:}
Từ điều kiện của $W$, ta viết lại hệ phương trình tuyến tính thuần nhất dưới dạng ma trận hóa như sau:
\[
\begin{pmatrix} 
1 & -1 & 2 & -1 & 1 \mid 0 \\ 
2 & -1 & 3 & -2 & 1 \mid 0 
\end{pmatrix}
\xrightarrow{(2) \to (2) - 2(1)}
\begin{pmatrix} 
1 & -1 & 2 & -1 & 1 \mid 0 \\ 
0 & 1 & -1 & 0 & -1 \mid 0 
\end{pmatrix}
\]
Do cột 3, 4, 5 không bán chuẩn hóa được nên hệ có vô số nghiệm như sau:
Đặt $\begin{cases} x_3 = a \\ x_4 = b \\ x_5 = c \end{cases} \quad \forall a, b, c \in \R$.
Từ dòng (2) ta có: $x_2 - x_3 - x_5 = 0 \Rightarrow x_2 = x_3 + x_5 = a + c$.
Từ dòng (1) ta có: $x_1 - x_2 + 2x_3 - x_4 + x_5 = 0 \Rightarrow x_1 = x_2 - 2x_3 + x_4 - x_5 = (a+c) - 2a + b - c = -a + b$.
Thay vào $W$ ta được:
\begin{align*}
W &= \{ X = (-a+b; a+c; a; b; c) \mid a, b, c \in \R \} \\
  &= \{ X = a(-1; 1; 1; 0; 0) + b(1; 0; 0; 1; 0) + c(0; 1; 0; 0; 1) \mid a, b, c \in \R \}.
\end{align*}
Đặt $\alpha_1 = (-1; 1; 1; 0; 0), \alpha_2 = (1; 0; 0; 1; 0), \alpha_3 = (0; 1; 0; 0; 1)$ và gọi $S = \{ \alpha_1, \alpha_2, \alpha_3 \}$.
$\Rightarrow W = \{ X = a\alpha_1 + b\alpha_2 + c\alpha_3 \mid a, b, c \in \R; \alpha_1, \alpha_2, \alpha_3 \in S \} = \langle S \rangle$.
Mà $\langle S \rangle \le V$ nên $W = \langle S \rangle \le V$ (đpcm).
Ta có $W = \langle S \rangle$ nên $S = \{ \alpha_1, \alpha_2, \alpha_3 \}$ là một hệ sinh của $W$.

Tiếp theo, ta tìm cơ sở và số chiều cho $W$ như sau:
Lập ma trận
\[
A_S = \begin{pmatrix} \alpha_1 \\ \alpha_2 \\ \alpha_3 \end{pmatrix} = \begin{pmatrix} -1 & 1 & 1 & 0 & 0 \\ 1 & 0 & 0 & 1 & 0 \\ 0 & 1 & 0 & 0 & 1 \end{pmatrix} 
\xrightarrow{(2) \to (2) + (1)} \begin{pmatrix} -1 & 1 & 1 & 0 & 0 \\ 0 & 1 & 1 & 1 & 0 \\ 0 & 1 & 0 & 0 & 1 \end{pmatrix} 
\]
\[
\xrightarrow{(3) \to (3) - (2)} \begin{pmatrix} -1 & 1 & 1 & 0 & 0 \\ 0 & 1 & 1 & 1 & 0 \\ 0 & 0 & -1 & -1 & 1 \end{pmatrix}
\]
Do ma trận đáp số có 3 dòng khác zero nên $r(A_S) = 3 = $ số véc tơ trong $S$, nên $S$ vừa là hệ sinh, vừa ĐLTT nên $S$ cũng là 1 cơ sở của $W$, và số chiều là $\dim_{\R} W = 3$.

\vspace{0.5cm}

\underline{Câu 2.} Trên $V=\R^3$, cho tập hợp
\[
a=\{u_1=(1; 1; 1); u_2=(1; 1; 0); u_3=(1; 0; 0)\}.
\]
\begin{enumerate}[label=\alph*/]
\item Chứng minh rằng $a$ là một cơ sở của $\R^3$.
\item Trực chuẩn hóa cơ sở $a$ thành cơ sở trực chuẩn $b$.
\item Tìm tọa độ của một véc tơ $v=(2; 3; 4) \in \R^3$ theo cơ sở $a$.
\item Tìm tọa độ của véc tơ $v=(2; 3; 4)$ theo cơ sở $b$.
\end{enumerate}

\underline{Giải:}
a/ Do $\R^3$ là không gian tuyến tính 3 chiều trên trường số thực $\R$, nên mỗi cơ sở của $\R^3$ là một hệ sinh, có 3 véc tơ, ĐLTT. Mà tập hợp $a$ đã có 3 véc tơ nên $a$ là 1 hệ sinh của $\R^3$. Để chứng tỏ $a$ là một cơ sở của $\R^3$ thì ta chỉ cần chứng tỏ $a$ là ĐLTT.
Ta lập ma trận
\[
A_a = \begin{pmatrix} u_1 \\ u_2 \\ u_3 \end{pmatrix} = \begin{pmatrix} 1 & 1 & 1 \\ 1 & 1 & 0 \\ 1 & 0 & 0 \end{pmatrix}
\]
Ta có $\det(A_a) = 1(0 - 0) - 1(0 - 0) + 1(0 - 1) = -1 \neq 0$.
$\Rightarrow a$ là một hệ sinh ĐLTT, nên $a$ là một cơ sở của $\R^3$.

b/ Dùng pp Gram--Schmidt, ta làm như sau:
\[
v_1 = u_1 = (1, 1, 1),
\]
\begin{align*}
v_2 &= u_2 - \frac{\ip{u_2}{v_1}}{\norm{v_1}^2}v_1 \\
    &= (1, 1, 0) - \frac{1(1) + 1(1) + 0(1)}{1^2 + 1^2 + 1^2}(1, 1, 1) \\
    &= (1, 1, 0) - \frac{2}{3}(1, 1, 1) = \left(\frac{1}{3}, \frac{1}{3}, -\frac{2}{3}\right),
\end{align*}
\begin{align*}
v_3 &= u_3 - \frac{\ip{u_3}{v_2}}{\norm{v_2}^2}v_2 - \frac{\ip{u_3}{v_1}}{\norm{v_1}^2}v_1 \\
    &= (1, 0, 0) - \frac{1(1/3) + 0(1/3) + 0(-2/3)}{(1/3)^2 + (1/3)^2 + (-2/3)^2}\left(\frac{1}{3}, \frac{1}{3}, -\frac{2}{3}\right) - \frac{1(1) + 0(1) + 0(1)}{3}(1, 1, 1) \\
    &= (1, 0, 0) - \frac{1/3}{6/9}\left(\frac{1}{3}, \frac{1}{3}, -\frac{2}{3}\right) - \frac{1}{3}(1, 1, 1) \\
    &= (1, 0, 0) - \frac{1}{2}\left(\frac{1}{3}, \frac{1}{3}, -\frac{2}{3}\right) - \left(\frac{1}{3}, \frac{1}{3}, \frac{1}{3}\right) = \left(\frac{1}{2}, -\frac{1}{2}, 0\right).
\end{align*}
Ta trực chuẩn hóa tiếp theo thành tập hợp trực chuẩn $b$:
\[
q_1 = \frac{v_1}{\norm{v_1}} = \frac{1}{\sqrt{3}}(1, 1, 1) = \left(\frac{1}{\sqrt{3}}, \frac{1}{\sqrt{3}}, \frac{1}{\sqrt{3}}\right),
\]
\[
q_2 = \frac{v_2}{\norm{v_2}} = \frac{1}{\sqrt{6/9}}\left(\frac{1}{3}, \frac{1}{3}, -\frac{2}{3}\right) = \frac{3}{\sqrt{6}}\left(\frac{1}{3}, \frac{1}{3}, -\frac{2}{3}\right) = \left(\frac{1}{\sqrt{6}}, \frac{1}{\sqrt{6}}, -\frac{2}{\sqrt{6}}\right),
\]
\[
q_3 = \frac{v_3}{\norm{v_3}} = \frac{1}{\sqrt{2/4}}\left(\frac{1}{2}, -\frac{1}{2}, 0\right) = \frac{2}{\sqrt{2}}\left(\frac{1}{2}, -\frac{1}{2}, 0\right) = \left(\frac{1}{\sqrt{2}}, -\frac{1}{\sqrt{2}}, 0\right).
\]
Suy ra $b = \{q_1, q_2, q_3\}$ là một cơ sở trực chuẩn cần tìm.

c/ Tìm tọa độ của $v=(2; 3; 4)$ theo cơ sở $a$.
Ta cần tìm $c_1, c_2, c_3 \in \R$ sao cho: $v = c_1u_1 + c_2u_2 + c_3u_3$
\begin{align*}
(2, 3, 4) &= c_1(1, 1, 1) + c_2(1, 1, 0) + c_3(1, 0, 0) \\
&\Leftrightarrow \begin{cases} 
c_1 + c_2 + c_3 = 2 \\ 
c_1 + c_2 = 3 \\ 
c_1 = 4 
\end{cases} \Leftrightarrow \begin{cases} 
c_1 = 4 \\ 
c_2 = -1 \\ 
c_3 = -1 
\end{cases}
\end{align*}
Như vậy, ta có: $[v]_a = \begin{pmatrix} c_1 \\ c_2 \\ c_3 \end{pmatrix} = \begin{pmatrix} 4 \\ -1 \\ -1 \end{pmatrix}$.

d/ Tìm tọa độ của $v=(2; 3; 4)$ theo cơ sở trực chuẩn $b$.
Do $b$ là cơ sở trực chuẩn của $\R^3$, nên ta có:
\[
d_1 = \ip{v}{q_1} = 2\left(\frac{1}{\sqrt{3}}\right) + 3\left(\frac{1}{\sqrt{3}}\right) + 4\left(\frac{1}{\sqrt{3}}\right) = \frac{9}{\sqrt{3}} = 3\sqrt{3},
\]
\[
d_2 = \ip{v}{q_2} = 2\left(\frac{1}{\sqrt{6}}\right) + 3\left(\frac{1}{\sqrt{6}}\right) + 4\left(-\frac{2}{\sqrt{6}}\right) = -\frac{3}{\sqrt{6}},
\]
\[
d_3 = \ip{v}{q_3} = 2\left(\frac{1}{\sqrt{2}}\right) + 3\left(-\frac{1}{\sqrt{2}}\right) + 4(0) = -\frac{1}{\sqrt{2}}.
\]
Suy ra $[v]_b = \begin{pmatrix} 3\sqrt{3} \\ -3/\sqrt{6} \\ -1/\sqrt{2} \end{pmatrix}$.

\vspace{0.5cm}

\underline{Câu 3:} Trên $V=\R^4$, với tích vô hướng $\ip{x}{y} = x_1y_1+x_2y_2+x_3y_3+x_4y_4$ (tích vô hướng tiêu chuẩn), cho hệ
$S=\{u_1=(1, -1, 1, -1); u_2=(1, 1, 1, 1); u_3=(1, 0, 0, 0)\}$.
Dùng pp Gram--Schmidt để trực chuẩn hóa $S$. 
Cho véc tơ $\lambda=(2, 1, 4, 1)$. Tìm tọa độ của véc tơ $\lambda$ theo cơ sở vừa trực chuẩn hóa được.

\underline{Giải:}
Dùng pp Gram--Schmidt, ta làm như sau:
\[
v_1 = u_1 = (1, -1, 1, -1),
\]
\begin{align*}
v_2 &= u_2 - \frac{\ip{u_2}{v_1}}{\norm{v_1}^2}v_1 \\
    &= (1, 1, 1, 1) - \frac{1(1) + 1(-1) + 1(1) + 1(-1)}{4}(1, -1, 1, -1) \\
    &= (1, 1, 1, 1) - 0(1, -1, 1, -1) = (1, 1, 1, 1),
\end{align*}
\begin{align*}
v_3 &= u_3 - \frac{\ip{u_3}{v_2}}{\norm{v_2}^2}v_2 - \frac{\ip{u_3}{v_1}}{\norm{v_1}^2}v_1 \\
    &= (1, 0, 0, 0) - \frac{1(1)}{4}(1, 1, 1, 1) - \frac{1(1)}{4}(1, -1, 1, -1) \\
    &= (1, 0, 0, 0) - \left(\frac{1}{4}, \frac{1}{4}, \frac{1}{4}, \frac{1}{4}\right) - \left(\frac{1}{4}, -\frac{1}{4}, \frac{1}{4}, -\frac{1}{4}\right) \\
    &= \left(\frac{1}{2}, 0, -\frac{1}{2}, 0\right).
\end{align*}
Ta trực chuẩn hóa tiếp theo:
\[
q_1 = \frac{v_1}{\norm{v_1}} = \frac{1}{\sqrt{4}}(1, -1, 1, -1) = \left(\frac{1}{2}, -\frac{1}{2}, \frac{1}{2}, -\frac{1}{2}\right),
\]
\[
q_2 = \frac{v_2}{\norm{v_2}} = \frac{1}{\sqrt{4}}(1, 1, 1, 1) = \left(\frac{1}{2}, \frac{1}{2}, \frac{1}{2}, \frac{1}{2}\right),
\]
\[
q_3 = \frac{v_3}{\norm{v_3}} = \frac{1}{\sqrt{1/4 + 1/4}}\left(\frac{1}{2}, 0, -\frac{1}{2}, 0\right) = \sqrt{2}\left(\frac{1}{2}, 0, -\frac{1}{2}, 0\right) = \left(\frac{1}{\sqrt{2}}, 0, -\frac{1}{\sqrt{2}}, 0\right).
\]
Tập hợp $b = \{q_1, q_2, q_3\}$ là cơ sở trực chuẩn vừa tìm được.

Vì $b$ là cơ sở trực chuẩn, tọa độ của $\lambda = (2, 1, 4, 1)$ theo $b$ được tính như sau:
\[
c_1 = \ip{\lambda}{q_1} = 2\left(\frac{1}{2}\right) + 1\left(-\frac{1}{2}\right) + 4\left(\frac{1}{2}\right) + 1\left(-\frac{1}{2}\right) = 1 - \frac{1}{2} + 2 - \frac{1}{2} = 2,
\]
\[
c_2 = \ip{\lambda}{q_2} = 2\left(\frac{1}{2}\right) + 1\left(\frac{1}{2}\right) + 4\left(\frac{1}{2}\right) + 1\left(\frac{1}{2}\right) = 1 + \frac{1}{2} + 2 + \frac{1}{2} = 4,
\]
\[
c_3 = \ip{\lambda}{q_3} = 2\left(\frac{1}{\sqrt{2}}\right) + 1(0) + 4\left(-\frac{1}{\sqrt{2}}\right) + 1(0) = -\frac{2}{\sqrt{2}} = -\sqrt{2}.
\]
Suy ra $[\lambda]_b = \begin{pmatrix} 2 \\ 4 \\ -\sqrt{2} \end{pmatrix}$.

\vspace{0.5cm}

\underline{Câu 4.} Cho ma trận vuông thực, cấp $2$, $A = \begin{pmatrix} 4 & 2 \\ 3 & -1 \end{pmatrix}$. Hãy chéo hóa ma trận $A$, rồi tìm $A^m$, $\forall m\in\N$ và suy ra $A^{2026}$.

\underline{Giải:}
Ta có đa thức đặc trưng của $A$:
\[
p_A(\lambda) = \det(A - \lambda I_2) = \begin{vmatrix} 4 - \lambda & 2 \\ 3 & -1 - \lambda \end{vmatrix} = (4 - \lambda)(-1 - \lambda) - 6 = \lambda^2 - 3\lambda - 10
\]
Giải pt: $p_A(\lambda) = 0 \Leftrightarrow \lambda^2 - 3\lambda - 10 = 0 \Leftrightarrow \begin{bmatrix} \lambda = 5 \\ \lambda = -2 \end{bmatrix}$.
Do đó, $A$ có 2 trị riêng: $\lambda_1 = 5, \lambda_2 = -2$ và số mũ là $r_1=1, r_2=1$.

Tiếp theo, ta tìm cơ sở $a_j$ cho các không gian riêng:
Với $\lambda_1 = 5$:
$E_5 = \{ X \in \R^2 \mid (A - 5I_2)X = O \}$
Ta có dạng ma trận hóa:
\[
\begin{bmatrix} -1 & 2 \\ 3 & -6 \end{bmatrix} \xrightarrow{(2) \to (2) + 3(1)} \begin{bmatrix} -1 & 2 \\ 0 & 0 \end{bmatrix}
\]
Hệ phương trình tương đương: $-x_1 + 2x_2 = 0 \Rightarrow x_1 = 2x_2$.
Đặt $x_2 = a \ (\forall a \in \R) \Rightarrow x_1 = 2a$.
$E_5 = \{ X = (2a, a) \mid a \in \R \} = \{ X = a(2, 1) \mid a \in \R \}$.
Đặt $u_1 = (2, 1)$, ta có $a_1 = \{u_1\}$ là cơ sở của $E_5$ và $\dim_{\R} E_5 = 1 = r_1$. \hfill (1)

Với $\lambda_2 = -2$:
$E_{-2} = \{ X \in \R^2 \mid (A + 2I_2)X = O \}$
Ta có dạng ma trận hóa:
\[
\begin{bmatrix} 6 & 2 \\ 3 & 1 \end{bmatrix} \xrightarrow{(2) \to (2) - \frac{1}{2}(1)} \begin{bmatrix} 6 & 2 \\ 0 & 0 \end{bmatrix}
\]
Hệ phương trình tương đương: $6x_1 + 2x_2 = 0 \Rightarrow x_2 = -3x_1$.
Đặt $x_1 = b \ (\forall b \in \R) \Rightarrow x_2 = -3b$.
$E_{-2} = \{ X = (b, -3b) \mid b \in \R \} = \{ X = b(1, -3) \mid b \in \R \}$.
Đặt $u_2 = (1, -3)$, ta có $a_2 = \{u_2\}$ là cơ sở của $E_{-2}$ và $\dim_{\R} E_{-2} = 1 = r_2$. \hfill (2)

Từ (1), (2) $\Rightarrow A$ có thể chéo hóa được trên $\R$.
Gọi $T = (u_1 \quad u_2) = \begin{pmatrix} 2 & 1 \\ 1 & -3 \end{pmatrix}$.
Ta có: $D = T^{-1}AT = \begin{pmatrix} 5 & 0 \\ 0 & -2 \end{pmatrix}$. Đây là dạng chéo hóa cần tìm của $A$.

Tìm $A^m$:
Ta có $T^{-1} = \frac{1}{-6 - 1}\begin{pmatrix} -3 & -1 \\ -1 & 2 \end{pmatrix} = \frac{1}{7}\begin{pmatrix} 3 & 1 \\ 1 & -2 \end{pmatrix}$.
\[
A^m = T D^m T^{-1} = \begin{pmatrix} 2 & 1 \\ 1 & -3 \end{pmatrix} \begin{pmatrix} 5^m & 0 \\ 0 & (-2)^m \end{pmatrix} \frac{1}{7}\begin{pmatrix} 3 & 1 \\ 1 & -2 \end{pmatrix}
\]
\[
= \frac{1}{7} \begin{pmatrix} 2\cdot 5^m & (-2)^m \\ 5^m & -3(-2)^m \end{pmatrix} \begin{pmatrix} 3 & 1 \\ 1 & -2 \end{pmatrix} 
= \frac{1}{7} \begin{pmatrix} 6\cdot 5^m + (-2)^m & 2\cdot 5^m - 2(-2)^m \\ 3\cdot 5^m - 3(-2)^m & 5^m + 6(-2)^m \end{pmatrix}
\]

Suy ra $A^{2026}$ (với $m = 2026$, lưu ý $(-2)^{2026} = 2^{2026}$):
\[
A^{2026} = \frac{1}{7} \begin{pmatrix} 6\cdot 5^{2026} + 2^{2026} & 2\cdot 5^{2026} - 2\cdot 2^{2026} \\ 3\cdot 5^{2026} - 3\cdot 2^{2026} & 5^{2026} + 6\cdot 2^{2026} \end{pmatrix}
\]

\begin{center}
----------HẾT----------
\end{center}

\end{document}
